\section{Experiments and Results}
In this section, we provide a comprehensive evaluation of the proposed Cross-Modality Semantic Alignment Transformer (CMSAT). We describe the experimental setup, implementation details, and present a series of comparative studies and analyses to demonstrate the effectiveness of leveraging LLM-generated fault semantics for zero-shot fault diagnosis in rotating machinery.

\subsection{Experimental Setup}
To validate the generalization capability of CMSAT, we employ two widely recognized datasets in the field of PHM: the Case Western Reserve University (CWRU) bearing dataset and the South East University (SEU) gearbox dataset.

The CWRU dataset \cite{key_cwru_dataset} consists of vibration signals collected from a motor-driven mechanical system. We categorize the faults into four types: normal, ball fault (BF), inner race fault (IRF), and outer race fault (ORF), each with three different severity levels (0.007, 0.014, and 0.021 inches). For the zero-shot task, we designate the normal state and a subset of fault types as seen classes, while the remaining fault categories at specific damage levels are treated as unseen classes. Specifically, the model is trained on 7 classes and tested on 3 completely unseen categories to evaluate its inductive reasoning performance.

The SEU dataset \cite{key_seu_dataset} represents a more complex mechanical environment, containing vibration data from a gearbox under two different load conditions (20 Hz-0 V and 30 Hz-2 V). Fault types include gear wear, tooth breakage, and various bearing fault configurations. Similar to the CWRU setup, we perform a distribution shift where the semantic descriptions of unseen gear faults are provided at inference time without any corresponding vibration samples during training.

Data preprocessing is critical for transformer-based architectures. All raw vibration signals are normalized using Z-score normalization to ensure zero mean and unit variance. We utilize a sliding window technique to segment the continuous signals into samples of length 1024, with an overlap ratio of 50\% to augment the training data. For each sample, the temporal signal is transformed into the frequency domain using the Fast Fourier Transform (FFT) to extract robust spectral features, which serve as the raw input to the vibration encoder. The semantic descriptions are generated using GPT-4, focusing on a multi-dimensional description of the mechanical physics, including vibration frequency range, impulse characteristics, and energy distribution.

\subsection{Implementation Details}
The CMSAT architecture is implemented using the PyTorch framework. The vibration encoder consists of 4 transformer layers with 8 multi-head attention blocks each, following the standard architectural principles for industrial signal processing \cite{key_transformer_base}. The embedding dimension $d_{model}$ is set to 512. For the semantic encoder, we utilize a pre-trained BERT-base model to extract initial text embeddings, which are then projected into the same 512-dimensional latent space to maintain modality compatibility \cite{key_bert_embedding}.

Training is conducted on an NVIDIA RTX 3090 GPU. We use the Adam optimizer with an initial learning rate of $1 \times 10^{-4}$ and a weight decay of $1 \times 10^{-5}$. A cosine annealing scheduler is applied to gradually reduce the learning rate. The batch size is set to 64, and the total number of epochs is 200. The temperature parameter $\tau$ in the contrastive loss function is initialized at 0.07. Training typically converges within 45 minutes for the CWRU dataset and 75 minutes for the SEU dataset. The optimal model is selected based on the validation loss of the seen classes.

\subsection{Zero-Shot Recognition Performance}
We compare the proposed CMSAT with five state-of-the-art method categories: standard deep learning (WDCNN), generative zero-shot models (ZSL-GAN), and advanced multi-modal frameworks (Deep Multi-modal ZSL, Attribute-based ZSL, and Semantic Embedding VRNN). Table~\ref{tab:comparison} summarizes the classification accuracy on unseen fault categories.

CMSAT significantly outperforms conventional ZSL methods, achieving an average accuracy of 89.4\% on the CWRU unseen set, which is a 5.2\% improvement over the strongest baseline, ZSL-GAN. Unlike GAN-based methods that rely on generating synthetic samples in the feature space, CMSAT directly aligns the signal features with the LLM-generated semantic space. This prevents the "mode collapse" common in generative models and ensures that the learned prototypes are physically meaningful.

The Evaluation process is grounded in the Zero-Shot Utility Function, defined as follows:
\begin{equation}
U = \frac{2 \times Acc_{seen} \times Acc_{unseen}}{Acc_{seen} + Acc_{unseen}}
\end{equation}
where $Acc_{seen}$ represents the accuracy on seen categories and $Acc_{unseen}$ denotes the accuracy on unseen categories \cite{key_gzsl_utility}. In Generalized Zero-Shot Learning (GZSL), the harmonic mean $U$ provides a more balanced measure than raw accuracy, preventing a bias toward seen classes that are more abundant during training. Our CMSAT maintains an $Acc_{seen}$ of 96.2\% while achieving a significant boost in $Acc_{unseen}$, resulting in a utility score that surpasses existing baseline results \cite{key_baseline_comparison}. Detailed diagnosis reveals the model's proficiency in distinguishing between minute fault variations, such as identifying the differences between an inner race fault and a ball fault at identical severity levels, even when these specific combinations were excluded from the training phase. This is largely due to the rich linguistic descriptors provided by the LLM, which encapsulate the fundamental physical interactions unique to each fault mechanism \cite{key_llm_phys_grounding}.

\subsection{Ablation Study}
To rigorously investigate the individual contribution of each component within our framework, we conduct a series of controlled ablation experiments. These studies focus on two primary dimensions: the source and richness of the semantic information, and the structural refinement of the modality fusion mechanism.

The first experiment compares the effectiveness of LLM-generated semantics against traditional, manually defined binary attribute vectors—a standard in earlier ZSL literature where researchers would manually tag features like "periodic impulses" or "high-frequency resonance" \cite{key_attr_vs_llm}. As summarized in Table~\ref{tab:ablation}, the transition from fixed attribute vectors to dynamic, natural language descriptions generated by an LLM results in an absolute performance gain of approximately 8.5\% for unseen categories. This improvement stems from the LLM's capacity to synthesize a "tapestry" of interconnected technical facts, such as how a specific bearing geometry influences the modulation sidebands under varying load conditions. Such nuanced information is far more descriptive than any discrete attribute set and allows the transformer architecture to build highly discriminative class prototypes in the latent space.

The second experiment evaluates our proposed Cross-Attention Alignment (CAA) mechanism against more basic integration methods, specifically a simple concatenation of vibration and semantic features followed by a standard Multi-layer Perceptron (MLP) \cite{key_fusion_mech}. The CAA module demonstrates a superior ability to bridge the modality gap. By computing attention scores between the segmented vibration signals and the semantic tokens, the model effectively masks out non-informative noise and focuses exclusively on the spectral components that align with the physical description of the fault. This leads to a significantly more compact and geometrically separable cluster distribution in the joint embedding manifold, as the model "learns" which physical frequencies correspond to specific linguistic terms used by the LLM.

\subsection{Parameter Sensitivity and Robustness}
The performance of the CMSAT is further scrutinized across various hyperparameter configurations and environmental noise levels to ensure its robustness in industrial settings. We perform a sensitivity analysis on several key parameters, most notably the embedding dimension $d_{z}$ and the temperature parameter $\tau$ used in the contrastive loss function.

Analysis of original embedding dimensions shows that a dimension of 512 strikes the optimal balance between representational expressiveness and computational footprint \cite{key_hyperparam_opt}. Dimensions lower than 256 result in an information bottleneck where subtle fault signatures are compressed and lost, while dimensions exceeding 1024 lead to incremental overfitting on the seen class distributions, thereby reducing the zero-shot generalization to unseen types. Similarly, the temperature parameter $\tau$ is found to be most effective at a value of 0.07, facilitating a sharp separation between positive and negative semantic pairs in the contrastive space.

In many practical scenarios, vibration signals are captured in high-noise environments that can mask fault signatures. To test the robustness of our approach, we inject additive Gaussian white noise into the raw vibration data at Signal-to-Noise Ratios (SNR) ranging from -10 dB to 10 dB. As visualised in Fig.~\ref{fig:snr}, CMSAT displays remarkable resilience, maintaining an unseen classification accuracy above 76\% even at an SNR of -4 dB. This robustness is a direct consequence of the cross-modal alignment; even when the vibration signal is distorted, the semantic prototypes derived from the LLM act as stable reference points or "anchor points" that guide the noisy vibration features toward their correct categorical regions \cite{key_robustness_study}. This suggests that the inclusion of language-based semantic knowledge provides a form of "regularization" that is absent in purely signal-based diagnostic models.

\subsection{Visual Analysis}
To provide a more intuitive and qualitative understanding of the latent space organization, we utilize t-Distributed Stochastic Neighbor Embedding (t-SNE) for the visualization of high-dimensional feature distributions and their corresponding semantic prototypes \cite{key_visual_tsne}.

The resulting visualization in Fig.~\ref{fig:tsne} shows the projected vibration features for unseen fault classes on both the CWRU and SEU datasets. Each cluster of signal features is clearly gravitating toward its respective LLM-generated semantic prototype, with distinct margins separating different fault types. The absence of significant overlap between the "Inner Race Fault" and "Outer Race Fault" samples—which often present similar periodic characteristics—confirms that the model has successfully internalized the unique semantic descriptions provided during the grounding phase.

Furthermore, we extract and visualize the attention maps generated by the Cross-Attention blocks during the inference of a "Gear Tooth Breakage" samples. The maps reveal that the model concentrates its focus on specific frequency bands associated with the gear rotational frequency and its upper-order harmonics, exactly as described in the LLM's output \cite{key_attention_visualization}. This transparency in the alignment process provides a level of interpretability that is crucial for trust in automated diagnostic systems. By showing that the transformer "looks" at the right physical indicators to match the semantic description, we validate that the CMSAT is not merely fitting noise but is instead performing a physically-informed diagnostic task.
